% !TeX program = xelatex
% CurveLab - standalone vector export. Compile with XeLaTeX, not pdfLaTeX.
% Structure: tailieuvehinh.pdf, I.1 (p.4); coordinates, draw, node, scope, clip.
% 3D is the current parallel projection, with visible/hidden line parts (I.3.5, pp.30-32).
% Captures the current viewport and sampled paths, without embedding a PNG.
% Width 16 cm; line widths and label sizes preserve the current drawing proportions.
% Mode: Hinh hoc khong gian
\documentclass[varwidth=50cm,border=2pt]{standalone}
\usepackage{fontspec}
\IfFontExistsTF{Times New Roman}{\setmainfont{Times New Roman}}{\setmainfont{TeX Gyre Termes}}
\IfFontExistsTF{Arial}{\setsansfont{Arial}}{\setsansfont{TeX Gyre Heros}}
\usepackage{amsmath,amssymb}
\usepackage{tikz,graphicx}
\usetikzlibrary{calc,angles,arrows.meta}
\definecolor{cl0}{RGB}{120,59,176}
\definecolor{cl1}{RGB}{130,144,160}
\definecolor{cl2}{RGB}{20,99,214}
\definecolor{cl3}{RGB}{176,53,135}
\definecolor{cl4}{RGB}{8,126,98}
\definecolor{cl5}{RGB}{193,90,16}
\definecolor{cl6}{RGB}{255,255,255}
\definecolor{cl7}{RGB}{121,131,145}
% Change only this width; geometry, strokes and labels scale together.
\newcommand{\FigureWidth}{16cm}
\begin{document}
\begin{center}
\resizebox{\FigureWidth}{!}{%
\begin{tikzpicture}[x=1cm,y=-1cm,
 declare function={figwidth=16; figheight=9.333333;},
  s1/.style={draw=cl0,draw opacity=1,line width=0.853583pt,line cap=butt,line join=miter},
  s2/.style={draw=cl1,draw opacity=1,line width=0.569055pt,line cap=butt,line join=miter,dash pattern=on 2.371063pt off 1.89685pt},
  s3/.style={fill=cl2,fill opacity=1},
  s4/.style={draw=cl2,draw opacity=1,line width=0.616476pt,line cap=butt,line join=miter},
  s5/.style={anchor=base west,inner sep=0pt,outer sep=0pt,text=cl2,text opacity=1,font={\rmfamily\bfseries\fontsize{7.113189}{8.535827}\selectfont}},
  s6/.style={fill=cl3,fill opacity=1},
  s7/.style={draw=cl3,draw opacity=1,line width=0.616476pt,line cap=butt,line join=miter},
  s8/.style={anchor=base west,inner sep=0pt,outer sep=0pt,text=cl3,text opacity=1,font={\rmfamily\bfseries\fontsize{7.113189}{8.535827}\selectfont}},
  s9/.style={fill=cl4,fill opacity=1},
  s10/.style={draw=cl4,draw opacity=1,line width=0.616476pt,line cap=butt,line join=miter},
  s11/.style={anchor=base west,inner sep=0pt,outer sep=0pt,text=cl4,text opacity=1,font={\rmfamily\bfseries\fontsize{7.113189}{8.535827}\selectfont}},
  s12/.style={fill=cl5,fill opacity=1},
  s13/.style={draw=cl5,draw opacity=1,line width=0.616476pt,line cap=butt,line join=miter},
  s14/.style={anchor=base west,inner sep=0pt,outer sep=0pt,text=cl5,text opacity=1,font={\rmfamily\bfseries\fontsize{7.113189}{8.535827}\selectfont}},
  s15/.style={fill=cl6,fill opacity=1},
  s16/.style={anchor=base west,inner sep=0pt,outer sep=0pt,text=cl7,text opacity=1,font={\rmfamily\bfseries\fontsize{7.113189}{8.535827}\selectfont}}
]
% Coordinates use the displayed projection; original coordinates are retained in comments.
\path[use as bounding box] (0,0) rectangle ({figwidth},{figheight});
\begin{scope}
\clip (0,0) rectangle ({figwidth},{figheight});
% A; world = [0,0,0]
\coordinate (P1) at (7.032179,9.085229);
% B; world = [4,0,0]
\coordinate (P2) at (11.535279,6.360014);
% C; world = [0,3,0]
\coordinate (P3) at (4.464721,6.396596);
% S; world = [1,1,4]
\coordinate (P4) at (7.302135,4.084438);
% H; world = [1,1,0]
\coordinate (P5) at (7.302135,7.507714);
% Drawing in display order: axes/grid, shapes, labels, formula.
\draw[s1] (P4)--(P1) (P4)--(P2) (P4)--(P3) (P1)--(P2) (P1)--(P3);
\draw[s2] (P2)--(P3) (7.217931,8.972814)--(7.076721,8.824939) (7.076721,8.824939)--(6.890968,8.937355);
\fill[s3] (P1) ellipse [x radius=0.083333cm,y radius=0.083333cm];
\draw[s4] (P1) ellipse [x radius=0.083333cm,y radius=0.083333cm];
\node[s5] at (7.182179,8.935229) {A};
\fill[s6] (P2) ellipse [x radius=0.083333cm,y radius=0.083333cm];
\draw[s7] (P2) ellipse [x radius=0.083333cm,y radius=0.083333cm];
\node[s8] at (11.685279,6.210014) {B};
\fill[s9] (P3) ellipse [x radius=0.083333cm,y radius=0.083333cm];
\draw[s10] (P3) ellipse [x radius=0.083333cm,y radius=0.083333cm];
\node[s11] at (4.614721,6.246596) {C};
\fill[s12] (P4) ellipse [x radius=0.083333cm,y radius=0.083333cm];
\draw[s13] (P4) ellipse [x radius=0.083333cm,y radius=0.083333cm];
\node[s14] at (7.452135,3.934438) {S};
\draw[s2] (P4)--(P5) (7.302135,7.366504)--(7.487887,7.254089) (7.487887,7.254089)--(7.487887,7.395299);
\fill[s15] (P5) ellipse [x radius=0.083333cm,y radius=0.083333cm];
\draw[s4] (P5) ellipse [x radius=0.083333cm,y radius=0.083333cm];
\node[s16] at (7.452135,7.357714) {H};
\end{scope}
\end{tikzpicture}%
}
\end{center}
\end{document}
